% Generated draft: requires mathematical and source review.
\documentclass[11pt]{article}
\usepackage[margin=25mm]{geometry}
\usepackage{fontspec}
\setmainfont{Latin Modern Roman}
\usepackage{amsmath,amssymb,mathtools,unicode-math}
\setmathfont{Latin Modern Math}
\usepackage{enumitem,needspace,xurl,hyperref}
\setlength{\parindent}{0pt}
\setlength{\parskip}{5pt}
\setlength{\emergencystretch}{3em}
\newcommand{\N}{\mathbb{N}}
\newcommand{\Z}{\mathbb{Z}}
\newcommand{\Q}{\mathbb{Q}}
\newcommand{\R}{\mathbb{R}}
\newcommand{\C}{\mathbb{C}}
\newcommand{\F}{\mathbb{F}}
\newcommand{\A}{\mathbb{A}}
\newcommand{\PP}{\mathbb{P}}
\newcommand{\E}{\mathbb{E}}
\newcommand{\eps}{\varepsilon}
\newcommand{\dd}{\,\mathrm{d}}
\begin{document}
\section*{UnsolvedMath ID: 30005206}
\textbf{Source problem code:} catalogue entry
\subsection{Definitions and mathematical statement}
The supplied passage concerns incompressible three-dimensional Euler flows and mentions spatially periodic smooth solutions as background. The following gives a precise periodic interpretation; the source does not explicitly fix the spatial domain, the Hölder exponent, or whether Hölder regularity is spatial or joint in space and time.

Let \(\mathbb T^3=\mathbb R^3/\mathbb Z^3\), equipped with its flat distance \(d\) and volume measure \(\dd x\). For \(T>0\), let \(u\colon[0,T]\times\mathbb T^3\to\mathbb R^3\) be continuous and let \(p\in L^1_{\mathrm{loc}}((0,T)\times\mathbb T^3)\) be a pressure. They form a weak incompressible Euler flow if
\[
\partial_t u+\operatorname{div}_x(u\otimes u)+\nabla_x p=0,
\qquad \operatorname{div}_x u=0
\]
in distributions on \((0,T)\times\mathbb T^3\), where \((u\otimes u)_{ij}=u_i u_j\). Define the kinetic-energy density and total kinetic energy by
\[
e(t,x)=\frac12|u(t,x)|^2,
\qquad E(t)=\int_{\mathbb T^3}e(t,x)\,\dd x.
\]
The local energy inequality means
\[
\partial_t e+\operatorname{div}_x\bigl((e+p)u\bigr)\leq0
\]
in distributions. Equivalently, for every nonnegative \(\varphi\in C_c^\infty((0,T)\times\mathbb T^3)\),
\[
\int_0^T\!\int_{\mathbb T^3}
\left(e\,\partial_t\varphi+(e+p)u\cdot\nabla_x\varphi\right)\,\dd x\,\dd t\geq0.
\]
Interpret nontrivial total kinetic-energy dissipation as
\[
E(t)\leq E(s)\quad(0\leq s<t\leq T),
\qquad E(t_1)<E(t_0)\quad\text{for some }0\leq t_0<t_1\leq T.
\]
This interpretation requires a genuine energy loss, but does not prescribe an energy profile or require strict decrease on every time interval.

For \(0<\alpha<1\), one explicit spatial Hölder convention is
\[
\sup_{t\in[0,T]}\sup_{x\ne y}
\frac{|u(t,x)-u(t,y)|}{d(x,y)^\alpha}<\infty.
\]
Under this convention, a precise periodic reading of the extracted question is: do there exist \(T>0\), \(\alpha\in(0,1)\), and \((u,p)\) satisfying the weak Euler equations, this Hölder condition, the local energy inequality, and nontrivial total kinetic-energy dissipation?

If joint space-time Hölder regularity is intended, the displayed spatial condition must instead be replaced by
\[
\sup_{(t,x)\ne(s,y)}
\frac{|u(t,x)-u(s,y)|}{\bigl(|t-s|+d(x,y)\bigr)^\alpha}<\infty.
\]
The supplied passage does not select between these conventions. Its existential question also does not specify whether constructions are sought for some positive exponent or for every exponent below \(1/3\). The supplied dated literature assessment additionally identifies a critical one-third formulation as unresolved. A critical formulation would require \(\alpha=1/3\), with the intended Hölder convention specified; it is a different quantifier requirement from existence for an unspecified positive exponent. These qualifications remain unresolved and the critical formulation is not attributed to the quoted passage as an explicit statement.
\subsection{Short English statement}
The question attributed to Hyunju Kwon's discussion asks for a Hölder-continuous three-dimensional incompressible Euler flow that loses total kinetic energy while satisfying the local energy inequality. A periodic weak formulation is given above, with genuine energy loss on at least one time interval. The exponent, time regularity, precise domain, and intended critical \(1/3\) target require clarification. The supplied label is partially solved; its literature assessment is dated curation, not independent verification.
\subsection{Sources}
\begin{itemize}
\item Supplied local UnsolvedMath registry record, permanent ID \(30005206\), title \emph{Strong Onsager Conjecture for Euler Flows}, source problem code \texttt{catalogue entry}, in the Oberwolfach Reports Open Problems collection.
\item Supplied primary report citation: \emph{Calculus of Variations (2023)}, Oberwolfach Reports, pp.\ 2129--2193, DOI \texttt{10.4171/owr/2022/37}. \url{https://doi.org/10.4171/owr/2022/37}. The excerpt is headed ``Euler flows and strong Onsager conjecture'' and names Hyunju Kwon. The registry gives proposed year 2022 but no separate proposer attribution.
\item Supplied literature reference: De Lellis--Kwon, dissipative Euler flows. \url{https://doi.org/10.2140/apde.2022.15.2003}.
\item Supplied literature reference: Giri--Kwon--Novack, \(L^3\)-based strong Onsager theorem. \url{https://arxiv.org/abs/2305.18509}.
\item The original Oberwolfach report is also listed among the supplied literature sources: \url{https://doi.org/10.4171/owr/2022/37}.
\item Literature assessment checked at \texttt{2026-08-22}: the supplied curation labels the problem partially solved and reports known subcritical Hölder and near-critical Besov constructions satisfying the local energy inequality, while describing a critical one-third Hölder formulation as unresolved. This is supplied dated curation, not an independently verified assessment of the literature or its current status.
\item The supplied excerpt attributes the resolution of the Onsager conjecture to Isett and refers to reference [1], but supplies no bibliographic entry for [1]. No missing citation is reconstructed here.
\end{itemize}
\end{document}
